% Unit 7 free-response set -- 2 long (10) + 3 short (4) = 32 pts.
%
% The ICE question is deliberately chosen so the 5% approximation FAILS
% (x is 20.5% of the initial concentration) and the quadratic is required.
% The chapter-track Ch13 set uses a perfect-square case solvable by taking
% roots; this one tests the judgement call the chapter set cannot.
\documentclass{apchem-exam}

\apcunit{7}
\apcunittitle{Equilibrium}
\apcdoctitle{Free-Response Set}
\apcsubtitle{CED 7.1--7.12 \textbullet\ 5 questions \textbullet\ 32 points \textbullet\ 50 minutes}

\begin{document}
\apctitleblock

\begin{apcnote}[Scope --- one CED exclusion applies]
Covers \ced{7.1} through \ced{7.12}, worth
\SI{7}{\percent}--\SI{9}{\percent} of the exam. Note that \ced{7.11} and
\ced{7.12} reach forward into solubility equilibria (Zumdahl \S16.1),
which is why $K_{sp}$ appears in an equilibrium unit.

Topic \ced{7.3} excludes \textbf{converting between $K_c$ and $K_p$}. The
conceptual difference between them is examinable; the
$K_p = K_c(RT)^{\Delta n}$ conversion is not, and is not asked below.
\end{apcnote}

\examsection{II}{Free Response}{5 questions \textbullet\ 32 points \textbullet\ 50 minutes}{\calculatorok}

% =====================================================================
\frq{long}{10}{\ced{7.4} \ced{7.7} \textbullet\ an ICE table where the approximation fails}

At a certain temperature $K_c = \num{0.0211}$ for
\[ \ce{PCl5(g) <=> PCl3(g) + Cl2(g)} \]
A vessel is charged with \SI{0.500}{\Molar} \ce{PCl5} and nothing else.

\begin{parts}
  \item Write the equilibrium expression and construct an ICE table.
        \answerlines[3]{$K_c = \dfrac{[\ce{PCl3}][\ce{Cl2}]}
        {[\ce{PCl5}]}$

        Initial: \num{0.500}, 0, 0. \quad Change: $-x$, $+x$, $+x$. \quad
        Equilibrium: $0.500-x$, $x$, $x$.}
  \item Show that the approximation $0.500 - x \approx 0.500$ is
        \emph{not} valid here. \workspace[2.8cm]
        \keyonly{\begin{apcanswer}Test it first. Assuming it holds,
        $x \approx \sqrt{(0.0211)(0.500)} = \num{0.103}$.

        As a fraction of the initial concentration that is
        $0.103/0.500 = \SI{20.5}{\percent}$ --- far above the
        \SI{5}{\percent} threshold. Dropping $x$ would change the
        denominator by a fifth, so the approximation must be
        \textbf{rejected} and the quadratic solved
        properly.\end{apcanswer}}
  \item Calculate the equilibrium concentration of each species.
        \workspace[3.6cm]
        \keyonly{\begin{apcanswer}$\dfrac{x^2}{0.500-x} = 0.0211
        \;\Rightarrow\; x^2 + 0.0211x - 0.01055 = 0$

        $x = \dfrac{-0.0211 + \sqrt{(0.0211)^2 + 4(0.01055)}}{2}
        = \dfrac{-0.0211 + 0.2065}{2} = \num{0.0927}$

        $[\ce{PCl3}] = [\ce{Cl2}] = \mathbf{0.0927}$~M \quad
        $[\ce{PCl5}] = 0.500 - 0.0927 = \mathbf{0.407}$~M

        Check: $(0.0927)^2/0.407 = \num{0.0211}$
        \checkmark\end{apcanswer}}
  \item The volume of the vessel is doubled at constant temperature. State
        the direction of the shift and justify.
        \answerlines[2]{The equilibrium shifts \textbf{right}. Doubling
        the volume lowers every partial pressure, and the system responds
        by moving toward the side with \emph{more} moles of gas --- one on
        the left, two on the right.}
\end{parts}

\begin{rubric}
  \rpoint{2}{(a) 1 pt the expression; 1 pt a correct ICE table}
  \rpoint{3}{(b) 1 pt computing the approximate $x$; 1 pt expressing it as
    a percentage of the initial concentration; 1 pt comparing with the
    \SI{5}{\percent} rule and rejecting. A candidate who never tests the
    approximation earns 0 here even if (c) is right}
  \rpoint{3}{(c) 1 pt the quadratic set up correctly; 1 pt
    $x = \num{0.0927}$; 1 pt all three concentrations}
  \rpoint{2}{(d) 1 pt shifts right; 1 pt toward more moles of gas, with
    the 1-versus-2 count}
\end{rubric}

% =====================================================================
\frq{long}{10}{\ced{7.5} \ced{7.9} \ced{7.10} \ced{7.11} \ced{7.12} \textbullet\ Q, Le Ch\^{a}telier and solubility}

\begin{parts}
  \item For the same reaction ($K_c = \num{0.0211}$), a mixture is found
        to contain $[\ce{PCl5}] = \SI{0.10}{\Molar}$,
        $[\ce{PCl3}] = \SI{0.20}{\Molar}$ and
        $[\ce{Cl2}] = \SI{0.20}{\Molar}$. Determine the direction of net
        change. \workspace[2.8cm]
        \keyonly{\begin{apcanswer}$Q = \dfrac{(0.20)(0.20)}{0.10}
        = \mathbf{0.40}$

        $Q > K$ ($0.40$ against $0.0211$), so there is too much product.
        The system shifts \textbf{left}, consuming \ce{PCl3} and \ce{Cl2},
        until $Q$ falls to $K$.\end{apcanswer}}
  \item State what the magnitude of $K_c = \num{0.0211}$ tells you about
        the equilibrium mixture. \answerlines[2]{$K \ll 1$, so at
        equilibrium the mixture is dominated by \textbf{reactant}: most of
        the \ce{PCl5} remains undecomposed.}
  \item A student claims that because $K$ is small, the reaction must be
        slow. Evaluate. \answerlines[3]{The claim is \textbf{wrong}. $K$
        describes only the \emph{position} of equilibrium --- how far the
        reaction proceeds --- and says nothing about how quickly it gets
        there. Rate is governed by activation energy, an entirely separate
        quantity. A reaction with a tiny $K$ may reach that equilibrium
        almost instantly.}
  \item Calcium fluoride has $K_{sp} = \num{3.9e-11}$. Calculate its molar
        solubility in pure water, then state and justify the effect of
        dissolving it in \SI{0.10}{\Molar} \ce{NaF} instead.
        \workspace[3.4cm]
        \keyonly{\begin{apcanswer}\ce{CaF2(s) <=> Ca^2+ + 2F^-}, so
        $[\ce{Ca^2+}] = s$ and $[\ce{F-}] = \mathbf{2}s$:

        $K_{sp} = (s)(2s)^2 = 4s^3 = 3.9\times10^{-11}$

        $s^3 = 9.75\times10^{-12} \;\Rightarrow\;
        s = \mathbf{2.1\times10^{-4}}$~M

        In \SI{0.10}{\Molar} \ce{NaF} the solubility \textbf{falls
        sharply}. Fluoride is a \textbf{common ion} and a product of the
        dissolution equilibrium, so adding it shifts the equilibrium left
        and less solid dissolves. $K_{sp}$ itself is unchanged --- it
        depends only on temperature.\end{apcanswer}}
\end{parts}

\begin{rubric}
  \rpoint{2}{(a) 1 pt $Q = \num{0.40}$; 1 pt shifts left because $Q > K$}
  \rpoint{2}{(b) 1 pt $K \ll 1$; 1 pt reactant-dominated}
  \rpoint{2}{(c) 1 pt the claim is wrong; 1 pt $K$ gives extent, not rate}
  \rpoint{4}{(d) 1 pt $[\ce{F-}] = 2s$; 1 pt the $4s^3$ form and
    $s = \SI{2.1e-4}{\Molar}$; 1 pt solubility decreases; 1 pt the
    common-ion shift with $K_{sp}$ unchanged. Writing $K_{sp} = s^3$ or
    claiming $K_{sp}$ falls each costs a point}
\end{rubric}

% =====================================================================
\frq{short}{4}{\ced{7.6} \textbullet\ heterogeneous equilibria}

\begin{parts}
  \item Write the equilibrium expression for
        \ce{NH4HS(s) <=> NH3(g) + H2S(g)}.
        \answerlines[2]{$K = [\ce{NH3}][\ce{H2S}]$ --- the solid is
        omitted.}
  \item Adding more solid \ce{NH4HS} to the sealed container does not
        change the pressure of \ce{NH3}. Explain.
        \answerlines[3]{A pure solid has a fixed density and composition,
        so its ``concentration'' does not change with how much is present.
        It therefore does not appear in $K$, and adding more cannot shift
        the equilibrium. As long as \emph{some} solid is present, the gas
        pressures are fixed by $K$ alone.}
\end{parts}

\begin{rubric}
  \rpoint{2}{(a) both gases included, solid omitted; including the solid
    earns 0}
  \rpoint{2}{(b) 1 pt a pure solid has constant concentration; 1 pt so its
    amount cannot shift the equilibrium}
\end{rubric}

% =====================================================================
\frq{short}{4}{\ced{7.1} \ced{7.2} \ced{7.8} \textbullet\ the dynamic nature of equilibrium}

\begin{parts}
  \item A sealed flask holds \ce{H2O(l)} in equilibrium with
        \ce{H2O(g)}. A few molecules of the liquid are replaced with
        water containing the heavy isotope \ce{^{18}O}. After some time
        the heavy isotope is found in \emph{both} phases. Explain what
        this shows. \answerlines[3]{It shows the equilibrium is
        \textbf{dynamic}. Evaporation and condensation have not stopped;
        they continue at equal rates, so molecules keep crossing between
        phases even though the amount of liquid and vapour stays constant.
        If the process had genuinely halted, the labelled molecules would
        have stayed where they were put.}
  \item State what would be observed macroscopically during this time.
        \answerlines[2]{Nothing changes: the volume of liquid, the vapour
        pressure and the temperature all stay constant. The macroscopic
        view alone cannot distinguish a dynamic equilibrium from a stopped
        reaction, which is exactly why the isotope experiment matters.}
\end{parts}

\begin{rubric}
  \rpoint{2}{(a) 1 pt both processes continue; 1 pt at equal rates}
  \rpoint{2}{(b) 1 pt no observable change; 1 pt recognizing the
    macroscopic view cannot by itself prove the process is dynamic}
\end{rubric}

% =====================================================================
\frq{short}{4}{\ced{7.9} \ced{7.10} \textbullet\ Le Ch\^{a}telier with a temperature change}

For the exothermic equilibrium
\ce{2SO2(g) + O2(g) <=> 2SO3(g)}, $\Delta H < 0$:

\begin{parts}
  \item State and justify the effect of raising the temperature on the
        amount of \ce{SO3} and on the value of $K$.
        \answerlines[3]{The amount of \ce{SO3} \textbf{decreases}. The
        forward reaction is exothermic, so heat behaves as a product;
        adding heat shifts the system in the endothermic direction, to the
        left.

        $K$ also \textbf{decreases}. Temperature is the only change that
        alters the value of $K$ itself rather than merely shifting the
        position along a fixed $K$.}
  \item State the effect of adding a catalyst on the amount of \ce{SO3}.
        \answerlines[2]{\textbf{No change.} A catalyst speeds the forward
        and reverse reactions equally, so equilibrium arrives sooner but
        lies in the same place.}
\end{parts}

\begin{rubric}
  \rpoint{3}{(a) 1 pt \ce{SO3} decreases; 1 pt heat treated as a product
    for an exothermic reaction; 1 pt $K$ itself changes with temperature}
  \rpoint{1}{(b) no change, with the reason}
\end{rubric}

\examtotal

\end{document}
